\documentclass[11pt,a4paper]{article}

%================================================
% ENCODAGE ET LANGUE
%================================================

\usepackage[utf8]{inputenc}
\usepackage[T1]{fontenc}
\usepackage[french]{babel}

%================================================
% MISE EN PAGE
%================================================

\usepackage{geometry}
\geometry{margin=1.55cm}

%================================================
% PACKAGE GTOLI
%================================================

\usepackage{../styles/gtoli}

%================================================
% INFORMATIONS DU DOCUMENT
%================================================

\title{\textbf{Cours -- 4e secondaire}\\[0.2cm]
\large Titre du chapitre}

\author{GToli}
\date{\today}

%================================================
% DOCUMENT
%================================================

\begin{document}

\maketitle

%================================================
% OBJECTIFS
%================================================

\begin{cadreobjectif}
À la fin de ce chapitre, l’élève doit être capable de :
\begin{itemize}
    \item identifier la notion étudiée ;
    \item mobiliser correctement le vocabulaire scientifique ;
    \item appliquer une méthode structurée ;
    \item justifier les différentes étapes du raisonnement ;
    \item résoudre des exercices similaires et des variantes.
\end{itemize}
\end{cadreobjectif}

%================================================
\section{Introduction et contexte}
%================================================

\begin{cadrebleu}{Introduction}
Cette section permet :
\begin{itemize}
    \item d’introduire la notion ;
    \item de contextualiser le chapitre ;
    \item de relier le contenu aux acquis antérieurs ;
    \item d’annoncer les objectifs du cours.
\end{itemize}
\end{cadrebleu}

%================================================
\section{Prérequis}
%================================================

\begin{cadregris}{Prérequis nécessaires}
Avant de commencer ce chapitre, il est conseillé de maîtriser :
\begin{itemize}
    \item notion 1 ;
    \item notion 2 ;
    \item notion 3.
\end{itemize}
\end{cadregris}

%================================================
\section{Rappel théorique}
%================================================

\begin{cadretheorie}
Cette section contient :
\begin{itemize}
    \item définitions ;
    \item propriétés ;
    \item théorèmes ;
    \item conventions ;
    \item formules importantes ;
    \item vocabulaire spécifique.
\end{itemize}

\medskip

\textbf{À retenir :}
la théorie doit être comprise avant l’application méthodologique.
\end{cadretheorie}

%================================================
\section{Méthode générale}
%================================================

\begin{cadremethode}
Pour résoudre ce type de problème :

\begin{enumerate}
    \item analyser l’énoncé ;
    \item identifier les données importantes ;
    \item choisir la méthode adaptée ;
    \item appliquer les propriétés nécessaires ;
    \item rédiger les étapes ;
    \item vérifier la cohérence du résultat.
\end{enumerate}
\end{cadremethode}

%================================================
\section{Exemple guidé}
%================================================

\begin{cadreexemple}
\textbf{Exemple :}

\medskip

Présenter ici un exercice modèle entièrement détaillé.

\medskip

\textbf{Analyse :}
\begin{itemize}
    \item Que cherche-t-on ?
    \item Quelles données possède-t-on ?
    \item Quelle méthode utiliser ?
\end{itemize}

\medskip

\textbf{Résolution :}

Présenter les étapes avec justification.
\end{cadreexemple}

%================================================
\section{Erreur fréquente}
%================================================

\begin{cadreattention}
Cette section permet de mettre en évidence :
\begin{itemize}
    \item une confusion fréquente ;
    \item une erreur de raisonnement ;
    \item un oubli classique ;
    \item une mauvaise utilisation d’une formule.
\end{itemize}
\end{cadreattention}

%================================================
\section{Exercice d’application}
%================================================

\begin{cadreenonce}
\textbf{Exercice :}

Rédiger ici un exercice similaire à l’exemple guidé.
\end{cadreenonce}

%================================================
\section{Aide méthodologique}
%================================================

\begin{cadreaide}
Pour résoudre cet exercice :

\begin{itemize}
    \item identifier les données utiles ;
    \item choisir la bonne formule ;
    \item organiser les calculs ;
    \item justifier les étapes importantes.
\end{itemize}
\end{cadreaide}

%================================================
\section{Solution détaillée}
%================================================

\begin{cadresolution}
La solution doit :
\begin{itemize}
    \item présenter clairement les étapes ;
    \item justifier les choix ;
    \item expliciter les calculs ;
    \item interpréter le résultat final.
\end{itemize}
\end{cadresolution}

%================================================
\section{Exercices progressifs}
%================================================

\begin{cadreactivite}
\begin{enumerate}
    \item exercice d’application directe ;
    \item exercice avec légère adaptation ;
    \item exercice nécessitant plusieurs étapes ;
    \item exercice de transfert ;
    \item exercice de synthèse.
\end{enumerate}
\end{cadreactivite}

%================================================
\section{Synthèse}
%================================================

\begin{cadresynthese}
\begin{itemize}
    \item notion essentielle ;
    \item formule importante ;
    \item méthode à retenir ;
    \item erreur à éviter ;
    \item stratégie de résolution.
\end{itemize}
\end{cadresynthese}

%================================================
\section{Auto-évaluation}
%================================================

\begin{cadregris}{Je fais le point}
\begin{itemize}
    \item[$\square$] Je comprends la théorie.
    \item[$\square$] Je sais utiliser la méthode.
    \item[$\square$] Je sais justifier mes étapes.
    \item[$\square$] Je sais refaire l’exemple seul.
    \item[$\square$] Je peux résoudre une variante.
\end{itemize}
\end{cadregris}

\end{document}